\documentclass[11pt]{article}
\usepackage[a4paper,margin=1in]{geometry}
\usepackage{amsmath,amssymb,amsthm,booktabs,hyperref}
\newtheorem{theorem}{Theorem}
\newtheorem{lemma}{Lemma}
\newtheorem{corollary}{Corollary}
\theoremstyle{definition}
\newtheorem{definition}{Definition}
\theoremstyle{remark}
\newtheorem{remark}{Remark}
\newcommand{\st}[2]{\bigl(#1 \mid #2\bigr)}
\newcommand{\too}{\;\Longrightarrow^{*}\;}
\title{Why the Claimed BB(7) Champion Halts:\\ the Induction, Written Out on Paper}
\author{Vishnu Vaidyanathan\\ \small Independent researcher\\ \small \texttt{vishnu@edgi.in}}
\date{12 November 2026 \\ \small working note; the halting part is now proved in Coq (\texttt{BB7Halts.v})}
\begin{document}
\maketitle

\begin{abstract}\sloppy
The machine \texttt{1RB0RG\_\allowbreak 1RC1RB\_\allowbreak 0RD1RZ\_\allowbreak 0LE0LD\_\allowbreak 0LF1LE\_\allowbreak 1RA1LF\_\allowbreak 1LC0RB}
(found by Jeff Ketchersid) is listed as the BB(7) champion, with the claim marked
as not independently verified. In earlier notes we proved in Coq that the machine
reaches a clean tape with 19 units, that one \emph{move} rule describes each
macro step for every gap width, and that a \emph{gate} rule halts the machine.
What was missing is a proof that the run uses only these rules until the end.
This note gives that proof on paper. It is an induction on the gap width. It also
gives the exact value of the last block as a ladder of functions $G_1,G_2,\dots$
The argument is checked by computer on every case small enough to compute.
The halting part has since been proved in Coq (\texttt{BB7Halts.v}, no axioms):
the machine halts from a blank tape. The exact value of the last block is still
a paper result.
\end{abstract}

\section{What is already proved}

A \emph{clean tape} is written
\[
\st{P}{(g_1,Z_1),(g_2,Z_2),\dots,(g_r,Z_r)} .
\]
It means: a block of $P$ ones, then $g_1$ zeros, then a block of $Z_1$ ones, then
$g_2$ zeros, and so on, with the head in state G at the left of the first block.
We call $P$ the \emph{working block}, each $Z_i$ a \emph{counter}, and each $g_i$
a \emph{gap}. A gap of $2j$ cells is said to have \emph{width} $j$.

Three facts are proved in \texttt{BB7Moves.v}.

\begin{description}
\item[Start (\texttt{start}).] From a blank tape, after 3{,}954 steps the tape is
$\st{16}{(4,2),\,(2,2)^{18},\,(2,1)}$. Here $(2,2)^{18}$ means eighteen copies.
So there are 19 counters equal to $2$ (the \emph{units}) and a last counter $1$.

\item[Move (\texttt{move}).] Let the first gap have width $j\ge 1$ and let
$m=3P+3-2j$. If $m\ge 1$ and $Z\ge 2$, then
\[
\underbrace{1^{P}\;0^{2j}\;1^{Z}\;R}_{\text{before}}
\quad\longrightarrow\quad
\underbrace{(1100)^{\,j-1}\;1^{N}\;00\;1^{Z-2}\;R}_{\text{after}},
\qquad N=3P+8-2j .
\]
$R$ is any rest of the tape. We say the move is \emph{in range} when $m\ge 1$.

\item[Gate (\texttt{halts\_gen}).] If the tape is $\st{P}{(2j,1)}$ with nothing
after the last $1$, and $m=3P+3-2j$ is odd and at least $1$, the machine halts.
\end{description}

\section{The move, read on clean tapes}

Reading the move rule in the clean-tape notation gives four cases. Let $T$ be
any list of (gap, counter) pairs and $N=3P+8-2j$.

\begin{center}\small
\begin{tabular}{@{}lll@{}}
\toprule
Case & Before & After \\
\midrule
\textbf{T} (triple), $j=1$, $Z\ge 4$ & $\st{P}{(2,Z),T}$ & $\st{3P+6}{(2,Z-2),T}$ \\
\textbf{T$'$} (counter ends), $j=1$, $Z=2$ & $\st{P}{(2,2),(g,X),T}$ & $\st{3P+6}{(g+2,X),T}$ \\
\textbf{S} (swap), $j\ge 2$, $Z\ge 4$ & $\st{P}{(2j,Z),T}$ & $\st{2}{(2,2)^{j-2},(2,N),(2,Z-2),T}$ \\
\textbf{S$'$} (counter ends), $j\ge 2$, $Z=2$ & $\st{P}{(2j,2),(g,X),T}$ & $\st{2}{(2,2)^{j-2},(2,N),(g+2,X),T}$ \\
\bottomrule
\end{tabular}
\end{center}

In words. With a single gap the working block triples. With a wider gap the
tripled block becomes a new counter $N$, a fresh working block of $2$ is planted,
and $j-2$ extra counters of $2$ are planted between them. When a counter is used
up, its two gaps join into one wider gap.

\begin{lemma}[Even numbers]\label{even}
If $P$ is even and every counter except possibly the last is even, the same is
true after any move.
\end{lemma}
\begin{proof}
$3P+6$ is even. $N=3P+8-2j$ is even. The planted values are $2$. $Z-2$ is even.
\end{proof}

After the start, $P=16$ and the counters are $2,\dots,2,1$. So from now on
\emph{$P$ is always even, every counter but the last is even, and the last
counter stays $1$ and stays last.}

\section{The ladder functions}

\begin{definition}\label{G}
For even $P$ define
\[
G_1(P)=3P+6,\qquad
G_j(P)=G_{j-1}^{\,(3P+8-2j)/2}\bigl(p_{j-1}\bigr)\quad (j\ge 2),
\]
where $f^{\,n}$ means ``apply $f$ $n$ times'', and the constants $p_k$ are
\[
p_1=2,\qquad p_{k}=G_{k-1}(p_{k-1})\quad(k\ge 2).
\]
\end{definition}

So $p_2=12$, $p_3=G_2(12)=G_1^{20}(2)=5\cdot 3^{20}-3$, and for example
$G_2(P)=G_1^{(3P+4)/2}(2)=5\cdot 3^{(3P+4)/2}-3$.
$G_1$ is linear, $G_2$ is an exponential, $G_3$ is a tower of exponentials, and
each further $G_j$ is one level higher.

\section{The consume lemma}

\begin{lemma}[Consume]\label{consume}
Let $1\le j\le 199$. Let $P$ and $Z$ be even, $Z\ge 2$, and let the first move be
in range ($3P+3-2j\ge 1$). Then for every next pair $(g,X)$ and every $T$,
\[
\st{P}{(2j,Z),\,(g,X),\,T}\;\too\;\st{G_j^{\,Z/2}(P)}{(g+2j,\,X),\,T},
\]
and every move used on the way is in range.
\end{lemma}

In words: a counter $Z$ behind a gap of width $j$ is used up completely; the
working block becomes $G_j$ applied $Z/2$ times; and the next gap grows by $2j$.

\begin{proof}
By induction on $j$.

\emph{$j=1$.} Apply case T $(Z-2)/2$ times, then case T$'$ once. Each
application replaces $P$ by $3P+6=G_1(P)$, and the last one joins the gaps to
$g+2$. Every move has gap width $1$ and $P\ge 2$, so $m=3P+1\ge 7$.

\emph{$j\ge 2$.} We show one \emph{round}:
\begin{equation}\label{round}
\st{P}{(2j,Z),\dots}\;\too\;
\begin{cases}
\st{G_j(P)}{(2j,Z-2),\dots} & Z\ge 4,\\[2pt]
\st{G_j(P)}{(g+2j,X),T} & Z=2.
\end{cases}
\end{equation}
Repeating the round $Z/2$ times gives the lemma.

The first move is case S or S$'$. It gives
\[
\st{2}{(2,2)^{j-2},\,(2,N),\,E,\dots},\qquad N=3P+8-2j,
\]
where $E=(2,Z-2)$ if $Z\ge 4$ and $E=(g+2,X)$ if $Z=2$.

\emph{The planted counters.} The working block is $p_1=2$ and it meets a
counter $2$ behind a gap of width $1$. By the lemma for width $1$ the block
becomes $G_1(p_1)=p_2$ and the next gap becomes $2+2=4$. Now $p_2$ meets a
counter $2$ behind a gap of width $2$; by the lemma for width $2$ the block
becomes $p_3$ and the next gap becomes $6$. Going on like this, the $k$-th
planted counter is met by $p_k$ behind a gap of width $k$, for $k=1,\dots,j-2$,
and after all of them the tape is
\[
\st{p_{j-1}}{(2(j-1),N),\,E,\dots} .
\]

\emph{The big counter.} By the lemma for width $j-1$ applied to the counter $N$,
the block becomes $G_{j-1}^{\,N/2}(p_{j-1})$, which is $G_j(P)$ by
Definition~\ref{G}, and the gap in front of $E$ grows by $2(j-1)$. If
$E=(2,Z-2)$ the new gap is $2j$. If $E=(g+2,X)$ the new gap is $g+2j$. This is
\eqref{round}.

\emph{Range.} We must check $3P'+3-2j'\ge 1$ for each working block $P'$ that
meets a gap of width $j'$:
\begin{itemize}
\item $p_1=2$ meets width $1$: $m=7$.
\item $p_2=12$ meets width $2$: $m=35$.
\item $p_k$ with $k\ge 3$ meets width $k$. By Lemma~\ref{big} below
$p_k\ge 132$, so $m\ge 399-2k\ge 1$ because $k\le 199$.
\item After a round, $G_j(P)$ meets width $j$ again. By Lemma~\ref{big},
$G_j(P)\ge 132$, so again $m\ge 1$.
\end{itemize}
Also $N=m+5\ge 6$, so $N$ is an even counter with $N/2\ge 3$.
\end{proof}

\begin{lemma}[Blocks are big enough]\label{big}
For $j\ge 2$, whenever the first move is in range, $G_j(P)\ge 132$.
Hence $p_k\ge 132$ for $k\ge 3$.
\end{lemma}
\begin{proof}
$G_j(P)=G_{j-1}^{\,n}(p_{j-1})$ with $n=N/2\ge 3$. For $j=2$ this is
$G_1^{\,n}(2)\ge G_1^{3}(2)=132$, since $G_1$ is increasing and $G_1(x)>x$.
For $j\ge 3$ the last of the $n$ applications is a value of $G_{j-1}$ on an
in-range block, which is at least $132$ by induction.
\end{proof}

\begin{remark}
Lemmas \ref{consume} and \ref{big} refer to each other. They are proved together
by one induction on $j$: the proof for width $j$ uses both lemmas only for
smaller widths, and Lemma~\ref{big} for width $j$ itself only after the first
round of width $j$ is done.
\end{remark}

\begin{remark}
The bound $j\le 199$ is only there because Lemma~\ref{big} uses the crude number
$132$. The blocks are in fact far larger. The real run needs widths up to $21$.
\end{remark}

\section{The theorem}

\begin{theorem}\label{main}
Let $1\le u\le 197$ and let $P_0\ge 2$ be even. Define
\[
Q_0=P_0,\qquad Q_k=G_{k+1}(Q_{k-1})\quad(k=1,\dots,u).
\]
Then the machine, started on the clean tape
$\st{P_0}{(4,2),\,(2,2)^{u-1},\,(2,1)}$, uses only in-range moves, reaches
\[
\st{Q_u}{(2(u+2),\,1)},
\]
and halts there.
\end{theorem}
\begin{proof}
The first unit is behind a gap of width $2$. By Lemma~\ref{consume} with $j=2$
and $Z=2$ it is used up, the block becomes $G_2(Q_0)=Q_1$, and the gap of the
second unit grows from $2$ to $6$, width $3$. By the lemma with $j=3$ the block
becomes $Q_2$ and the third unit's gap becomes $8$. In general the $k$-th unit
is met by $Q_{k-1}$ behind a gap of width $k+1$, and is used up by the lemma
with $j=k+1$. The first moves are in range: $Q_0\ge 2$ meets width $2$
($m=3P_0-1\ge 5$), and $Q_k\ge 132$ for $k\ge 1$.

After the last unit the tape is $\st{Q_u}{(2(u+2),1)}$. $Q_u$ is even by
Lemma~\ref{even}, so $m=3Q_u+3-2(u+2)$ is odd, and it is at least $1$ because
$Q_u\ge 132$. The gate rule applies and the machine halts.
\end{proof}

\begin{corollary}[The champion]\label{champ}
The machine halts from a blank tape. The widest gap in a move is of width $20$,
the halt is at width $21$, and the last working block is
\[
Q_{19}=G_{20}\bigl(G_{19}(\cdots G_3(G_2(16))\cdots)\bigr).
\]
\end{corollary}
\begin{proof}
The start rule gives the tape of Theorem~\ref{main} with $P_0=16$ and $u=19$.
\end{proof}

The first value is $Q_1=G_2(16)=5\cdot 3^{26}-3=12{,}709{,}329{,}141{,}642$.
The second, $Q_2=G_3(Q_1)$, is $G_2$ applied
$19{,}063{,}993{,}712{,}464$ times to $12$: a tower of that many
exponentials. We do not work out a closed bound here. The claimed bound of the
form $3\uparrow^{20}3$ has the same shape as a ladder with twenty levels, but
comparing the two is left for later.

\section{Computer checks}

All of these are in the repository, in \texttt{bb7/}.

\begin{center}
\begin{tabular}{@{}p{0.5\textwidth}p{0.42\textwidth}@{}}
\toprule
Check & Result \\
\midrule
Move and halt words against the machine, digit for digit, widths $1$ to $21$
(\texttt{check\_wide\_gaps.py}) & 63 of 63 and 63 of 63 match \\
Every shape reachable with $1$ to $12$ units (\texttt{check\_closure.py}) &
no shape without a rule; widest move $u+1$; one halt, at width $u+2$ \\
Lemma~\ref{consume} against the exact rules, ten small cases
(\texttt{check\_induction.py}) & 10 of 10 agree (values up to 23{,}478 digits) \\
Theorem~\ref{main} with $u=1$ against the real machine & $P_0=2$: halts at step
18{,}310{,}165 as predicted, last block $1212=G_2(2)$.\quad $P_0=4$: halts at step
13{,}318{,}711{,}129 as predicted, last block $32802=G_2(4)$ \\
$Q_1$ and the counter $3Q_1+2$ against the macro run & both equal \\
\bottomrule
\end{tabular}
\end{center}

No case with two or more units can be run to the end; the numbers are too large.

\section{What is proved and what is not}

\begin{itemize}
\item \textbf{Proved in Coq} (\texttt{BB7Moves.v}, \texttt{BB7Halts.v}, Coq 8.18,
no axioms): the start, the move for every width, the gate, the consume lemma for
widths $1$ to $199$ (as ``the counter is used up, the next gap grows by $2j$, and
the new block is even and large enough''), and
\begin{center}
\texttt{champion\_halts}: there is an $n$ such that the machine, started on a
blank tape, is in the halt state after $n$ steps.
\end{center}
Also \texttt{champion\_gate}: the run reaches one block, a gap of $42$ cells
(width $21$) and the last $1$.
\item \textbf{Proved on paper only}: the exact values, that is, that the new
block after a consume is $G_j^{\,Z/2}(P)$ and that the last block is
$G_{20}(\cdots G_2(16))$. The Coq proof shows the block exists, is even and is
at least $132$; it does not name it.
\item The Coq version of Theorem~\ref{main} (\texttt{units\_halt}) asks for
$P_0\ge 12$ instead of $P_0\ge 2$. The champion has $P_0=16$.
\item \textbf{Not done}: the halting time, the number of ones left, and a
comparison with the claimed bound.
\item We do not know whether this argument is already known to the discoverer
or to others.
\end{itemize}

\section{How the Coq proof is arranged}

\begin{enumerate}
\item Numbers are written by their halves: the working block is $2q$ ones and a
counter is $2h+2$ ones. Then ``even'' needs no side condition.
\item \texttt{consume1}: width $1$, by induction on $h$.
\item \texttt{planted}: a row of $t$ units behind a gap of width $a+1$ is used
up at widths $a+1,\dots,a+t$, given the consume lemma for those widths.
\item \texttt{round}: one round at width $j\ge 2$: the move, then
\texttt{planted} for the $j-2$ planted units, then the consume lemma at width
$j-1$ for the big counter.
\item \texttt{consume\_all}: all widths up to $199$, by induction on the width
with an inner induction on $h$.
\item \texttt{units\_to\_gate}: \texttt{planted} once more, now for the 19 real
units ($a=1$). Then \texttt{halts\_gen} gives \texttt{champion\_halts}.
\end{enumerate}
Every lemma carries a free rest-of-tape, as the move theorem does. That is what
lets the consume lemma be used inside itself.

\end{document}
